% Generated draft: requires mathematical and source review.
\documentclass[11pt]{article}
\usepackage[margin=25mm]{geometry}
\usepackage{fontspec}
\setmainfont{Latin Modern Roman}
\usepackage{amsmath,amssymb,mathtools,unicode-math}
\setmathfont{Latin Modern Math}
\usepackage{enumitem,needspace,xurl,hyperref}
\setlength{\parindent}{0pt}
\setlength{\parskip}{5pt}
\setlength{\emergencystretch}{3em}
\newcommand{\N}{\mathbb{N}}
\newcommand{\Z}{\mathbb{Z}}
\newcommand{\Q}{\mathbb{Q}}
\newcommand{\R}{\mathbb{R}}
\newcommand{\C}{\mathbb{C}}
\newcommand{\F}{\mathbb{F}}
\newcommand{\A}{\mathbb{A}}
\newcommand{\PP}{\mathbb{P}}
\newcommand{\E}{\mathbb{E}}
\newcommand{\eps}{\varepsilon}
\newcommand{\dd}{\,\mathrm{d}}
\begin{document}
\section*{UnsolvedMath ID: 6700031}
\textbf{Source problem code:} catalogue entry
\subsection{Definitions and mathematical statement}
The source asks for the sharp values of
\[
\operatorname{width}_{n-m}(Q),\qquad m\geq2,
\]
for the objects \(Q\) referred to only as ``these solids''. Both the family of solids and the definition of \(\operatorname{width}_{n-m}\) occur outside the supplied excerpt. Consequently, neither a particular width invariant nor the admissible metrics, dimensional index conventions, and normalizations can be fixed from this item. If \(n\) denotes the dimension of a solid, the subscript is \(n-m\); the excerpt does not explicitly make even that convention or give the upper range of \(m\).

The request is to determine \emph{sharp values}, rather than merely some bound. It is expressly qualified by Gromov's caveat ``unless I missed some paper''. The research summary suggests metric balls and scalar curvature bounds, but supplies no particular solids, radius, curvature normalization, or formula. Those suggestions therefore do not define an additional hypothesis or an exact optimization problem. A self-contained numerical formulation requires the missing source context; it cannot be obtained by selecting one of the different invariants commonly called width.
\subsection{Short English statement}
Determine the exact higher codimension widths of the solids discussed in the missing surrounding passage, for \(m\geq2\). The excerpt does not identify the solids or define its width invariant.
\subsection{Sources}
\begin{itemize}
\item Local UnsolvedMath registry, permanent ID 6700031, source code catalogue entry; Misha Gromov, 2017; updated 2026-08-06.
\item Gromov, \emph{101 Questions, Problems and Conjectures around Scalar Curvature} (2017), printed tag [?31], occurrence 32, PDF page 42: \url{https://www.ihes.fr/~gromov/wp-content/uploads/2018/08/101-problemsOct1-2017.pdf}. The recorded access date is 2026-07-29; extraction is normalized PDF text from a public author-hosted incomplete PDF. Current status and reuse terms require review.
\item The supplied research summary mentions width estimates, Kumura-type bounds, spheres, balls, and the two-sphere lemma without a precise further citation or quantitative statement. The adjacent local records 6700030 and 6700032 do not identify ``these solids'' or resolve the width convention.
\end{itemize}
\end{document}
